\scanpage{44}
% archive sheet number: 43
% The circled text at the head is called in after « Fil¹ est unique et ».
\[
\phi^{\handwritten{*}}(\mathrm{Fil}^1(M_{S_{\unsure{o|0}}}))=
\operatorname{Ker} F_{M_o}\ (=\operatorname{Im}V_{M_o})\handwritten{\subset}M_o^{(p)}\quad .
\]
D'a\deleted{\unsure{n}}utre part, on montre que $\phi$ a des propriétés de fidélité telles que la relation précédente détermine sans ambiguité le sous-Module localement facteur direct $\mathrm{Fil}^1$ de $M_{S_o}$.

On dira qu'un \inserted{F-V-}c\deleted{\unsure{ri}}ristal en modules loc. lib. $M_o$ sur le site cristallin de $S_o$ au des\deleted{\unsure{s}}sus de $\underline{F}_p$ est \uline{admissible} s'il existe un sous-module l\deleted{\unsure{a}}oc. \deleted{facteu} facteur direct $\mathrm{Fil}^1$ de $M_{oS}$ tel que l'on ait la relation précédente\unsure{;|,} \deleted{On} \inserted{$\mathrm{Fil}^1$ est unique et \deleted{déterm}\unsure{$\circ$|o}\inserted{détermine une filtration qu'on appellera la \uline{filt}. \uline{canonique} . \handwritten{On}}} dit de même qu'un \deleted{\unsure{Fil}}cristal de Dieudonné \inserted{$M$} sur $S_o$ est admissible si \deleted{\unsure{l}}sa restriction $M_o$ \deleted{sur} au site cristallin rel. à $\underline{F}_p$ est admissible. On voit donc que si $M$ est le cristal de Dieudonné associé à un groupe de BT sur $S_o$, $M$ est admissible, et la filtration de $M_{S_o}$ envisagée dans b) est la filtra- canonique définie précédemment.

% slanted
\margin{\handwritten{d) Passage à la limite : \unsure{prendre|pour des|formule} anneau $A$ \unsure{séparé} et complet pour top. \uline{\unsure{$p$-adique}}}}
% slanted
\margin{\handwritten{e) Théorie \unsure{:|à} isogénie près \uline{et} th. de Tate}}
% slanted
\margin{\handwritten{f) Dualité}}

\handwritten{2}) \uline{Théorie de déformations pour les schémas abéliens}.

\handwritten{a)} Dans cette théorie, il n'y a plus de nombre premier privilégié, par contre il faut travailler avec le site cristallin \uline{nilpotent}. On trouve un foncteur \inserted{contravariant} de la catégorie des schémas abéliens sur le schéma $S$ vers les \deleted{sites} cristaux sur le site cristallin nilpotent de $S$
\[
\text{Schémab}(S)\handwritten{\longrightarrow}\mathrm{Crisloclib}_{\mathrm{nilp}}(S)
\]
qui peut être défini par exemple comme
\[
\underline{\handwritten{\mathbb{D}}}^{*}:\ A\handwritten{\longmapsto}R^1 f_{A\mathrm{cris}*}(\deleted{A}\handwritten{O_{A_{\mathrm{cris}}}})
\tag*{\handwritten{(*)}}
\]
où on considère
\[
f_{\mathrm{cris}}:\quad A_{\mathrm{crisnilp}}\handwritten{\longrightarrow}S_{\mathrm{crisnilp}}
\]
induit par $S$. Par \deleted{\unsure{l}}une va\deleted{\unsure{i}}riante du théorème de changement de base, on voit que
\[
\mathbb{D}^*(A)_S=\deleted{\unsure{R^1 f_{A\mathrm{DR}*} A}}\underline{H}_{\mathrm{DR}}^1(A/S)\overset{\mathrm{dfn}}{=}R^1 f_{A*}(\handwritten{\Omega}_{A/S}^{\bullet})\quad .
\]
Pour calculer $\mathbb{D}^{\handwritten{*}}(A)_{S'}$, $S'$ épaississement à p.div. loc. nilpotentes de $S$, on aura \deleted{\unsure{p}}pour tout prolongement \deleted{canonique} de $A$ en un schéma abélien $A'$ sur $S'$
